\documentclass[12pt,a4paper]{article}
\usepackage[utf8]{inputenc}
\usepackage[T1]{fontenc}
\usepackage{lmodern}
\usepackage[margin=2.5cm]{geometry}
\usepackage{hyperref}
\usepackage{enumitem}
\usepackage{xcolor}
\usepackage{listings}
\usepackage{booktabs}
\usepackage{float}
\usepackage{amsmath}
\usepackage{siunitx}
\usepackage{longtable}
\usepackage{tikz}
\usepackage{caption}
\usepackage{subcaption}

% ----- code listing style -----
\definecolor{codebg}{rgb}{0.95,0.95,0.95}
\lstset{
	backgroundcolor=\color{codebg},
	basicstyle=\ttfamily\small,
	breaklines=true,
	frame=single,
	language=C,
	tabsize=2,
	showstringspaces=false,
	numbers=left,
	numberstyle=\tiny\color{gray},
	keywordstyle=\color{blue}\bfseries,
	commentstyle=\color{green!60!black},
	stringstyle=\color{red},
	captionpos=b,
	morekeywords={HAL_StatusTypeDef, I2C_HandleTypeDef, TIM_HandleTypeDef, UART_HandleTypeDef,
		GPIO_InitTypeDef, RCC_OscInitTypeDef, RCC_ClkInitTypeDef, RCC_PeriphCLKInitTypeDef,
		encoder_data_t, pid_gains_t, mpu6050_axes_t, mpu6050_offsets_t, motor_id_t}
}

% ----- helpers -----
\newcommand{\code}[1]{\lstinline[basicstyle=\ttfamily\small]{#1}}
\newcommand{\file}[1]{\texttt{#1}}
\newcommand{\reg}[1]{\texttt{#1}}
\newcommand{\pin}[1]{\texttt{#1}}
\newcommand{\var}[1]{\texttt{#1}}
\newcommand{\func}[1]{\texttt{#1}}
\newcommand{\isr}[1]{\texttt{#1}}

% ----- document -----
\begin{document}
	
	\title{\Huge \textbf{Self‑Balancing Robot Firmware:}\\[0.3em]
		\Large A Complete Walk‑through of the STM32F3 + Keyestudio Kit Implementation}
	\author{Prepared for Lab 11 -- Open‑Ended Project}
	\date{\today}
	\maketitle
	
	\tableofcontents
	\newpage
	
	% ================================================================
	\section{Introduction}
	\label{sec:intro}
	% ================================================================
	
	This document explains, in painstaking detail, every module and every significant line of code that powers a two‑wheeled self‑balancing robot.  
	The robot is built around an \textbf{STM32F3Discovery} board (STM32F303VCT6, clocked at \SI{48}{\mega\hertz}) and the \textbf{Keyestudio Self‑Balancing Car Kit V3}.  
	The firmware reads an MPU‑6050 inertial measurement unit, fuses the sensor data to obtain a stable tilt angle, runs a PID controller 200 times per second, and drives two DC motors to keep the robot upright.
	
	Because the project is meant to be studied and questioned in a viva‑voce examination, this report leaves nothing to guesswork.  
	It starts from the hardware wires and timers, moves through every C source file, and finishes with a discussion of creative features that go beyond the strict requirements of \textbf{Lab 11} (Open‑Ended Lab).  
	Even if you have never written a line of embedded C before, you should be able to follow the explanations, and after reading you will be equipped to answer detailed questions about the code.
	
	\bigskip
	\noindent\textbf{How to use this document:}
	\begin{itemize}
		\item First read Sections~\ref{sec:overview}--\ref{sec:arch} to get the big picture.
		\item Then dive into the file‑by‑file analysis in Section~\ref{sec:code} for an in‑depth understanding.
		\item Finally, check Section~\ref{sec:extra} to see what makes this project stand out.
	\end{itemize}
	
	% ================================================================
	\section{Project Overview}
	\label{sec:overview}
	% ================================================================
	The firmware is built with STM32Cube HAL libraries, compiled with ARM GCC, and flashed from VS Code.  
	A hardware timer (TIM6) fires at exactly \SI{200}{\hertz} and inside its interrupt service routine (ISR) the sensor is read, the angle is updated, the PID law is computed, and the motors are commanded.  
	This guarantees that the balancing reaction is never delayed by other software tasks.
	
	A separate UART telemetry stream sends live data -- angle, gyro rate, motor command, and wheel encoder speeds -- to a PC for plotting.  
	The transmission is designed to be \textit{non‑blocking}, meaning it does not stall the main control loop.
	
	The code is split into independent modules, each with its own \texttt{.c} and \texttt{.h} file:
	\begin{itemize}
		\item \texttt{mpu6050} -- I²C driver for the IMU
		\item \texttt{complementary} -- sensor‑fusion filter
		\item \texttt{pid} -- generic PID controller
		\item \texttt{motor} -- DC motor driver (PWM + direction)
		\item \texttt{encoder} -- wheel‑speed measurement via input capture
		\item \texttt{main} -- application entry point, system initialisation, telemetry, and interrupt glue
	\end{itemize}
	
	% ================================================================
	\section{Hardware Setup and Pin Configuration}
	\label{sec:hardware}
	% ================================================================
	
	\subsection{Microcontroller and Shield}
	The \textbf{STM32F3Discovery} board features an STM32F303VCT6 microcontroller with an Arm Cortex‑M4 core.  
	The Keyestudio shield provides two DC motors with integrated magnetic encoders, an L293D‑compatible H‑bridge, an on‑board MPU‑6050 module, and a 5\,V regulator for the encoder LEDs.
	
	\subsection{Pin Mapping}
	Table~\ref{tab:pins} lists every connection between the Discovery board and the shield.  
	\textbf{Common ground is essential:} the STM32 ground must be tied to the shield’s GND.
	
	\begin{table}[H]
		\centering
		\caption{Complete pin mapping -- verified against firmware \file{motor.h}, \file{encoder.c}, and \file{main.c}}
		\label{tab:pins}
		\begin{tabular}{@{}lll@{}}
			\toprule
			\textbf{STM32 Pin} & \textbf{Shield Function} & \textbf{Description} \\
			\midrule
			PA0  & TIM2\_CH1 & Right motor encoder input (Input Capture) \\
			PA1  & TIM2\_CH2 & Left motor encoder input (Input Capture) \\
			PB10 & TIM2\_CH3 & Left motor PWM output \\
			PB11 & TIM2\_CH4 & Right motor PWM output \\
			PD8  & D12 (Right DirA) & Direction control, right motor A \\
			PD9  & D8  (Right DirB) & Direction control, right motor B \\
			PD10 & D10 (Left DirA)  & Direction control, left motor A  \\
			PD11 & D7  (Left DirB)  & Direction control, left motor B  \\
			PA9  & I2C2\_SCL & MPU‑6050 serial clock \\
			PA10 & I2C2\_SDA & MPU‑6050 serial data \\
			PA2  & USART2\_TX & Telemetry UART transmit \\
			PA3  & USART2\_RX & Telemetry UART receive (optional) \\
			GND  & GND & Common ground (must be connected) \\
			5V   & 5\,V (shield encoder) & Powers encoder modules (max \SI{100}{\milli\ampere}) \\
			\bottomrule
		\end{tabular}
	\end{table}
	
	\noindent\textbf{Important safety note:} The motor supply (\texttt{Vin} on the shield) must \textbf{never} be connected to the STM32’s 5\,V or 3.3\,V pins. Use a separate battery capable of at least \SI{1}{\ampere}.
	
	\subsection{Timer and Peripheral Configuration}
	The firmware uses two hardware timers:
	
	\begin{itemize}
		\item \textbf{TIM2} -- a 16‑bit timer clocked by the APB1 bus (after prescaler).  
		\begin{itemize}
			\item Prescaler = 2, Auto‑reload = 1599 → PWM frequency = \(\frac{48\,\text{MHz}}{2+1} / (1599+1) = 10\,\text{kHz}\).
			\item Channels 1 \& 2: Input Capture (rising edge only), used for the encoder pulse measurement.  
			Channel~1 → Left motor encoder, Channel~2 → Right motor encoder.
			\item Channels 3 \& 4: PWM generation in mode 1.  
			Channel~3 → Left motor (PB10), Channel~4 → Right motor (PB11).
		\end{itemize}
		\item \textbf{TIM6} -- a basic timer used as the control‑loop heartbeat.
		\begin{itemize}
			\item Prescaler = 239, Auto‑reload = 999.
			\item Timer clock = \(\frac{48\,\text{MHz}}{(239+1)} = 200\,\text{kHz}\).
			\item Overflow frequency = \(\frac{200\,\text{kHz}}{(999+1)} = 200\,\text{Hz}\) → exactly \SI{5}{\milli\second} period.
		\end{itemize}
	\end{itemize}
	
	\noindent\textbf{I²C2} is configured in Fast Mode (standard \SI{400}{\kilo\hertz}) with analog and digital filters enabled.  
	\textbf{USART2} runs asynchronously at 115200\,baud, 8 data bits, no parity.
	
	% ================================================================
	\section{System Architecture and Control Loop}
	\label{sec:arch}
	% ================================================================
	
	The robot’s brain operates in two layers:
	\begin{enumerate}
		\item \textbf{Fast layer (TIM6 ISR)} -- every \SI{5}{\milli\second}: read sensor → fuse angle → compute PID → drive motors.
		\item \textbf{Slow layer (main loop)} -- every \SI{100}{\milli\second} (approximately): send telemetry and print a heartbeat.
	\end{enumerate}
	
	\subsection{Angle Estimation -- Complementary Filter}
	The tilt angle cannot be measured directly; it must be \textit{estimated}. The MPU‑6050 provides two clues:
	\begin{itemize}
		\item \textbf{Accelerometer} measures the direction of gravity. When the robot is still, the angle is \(\tan^{-1}(a_y / a_z)\). However, during movement the accelerometer also picks up linear acceleration, making it noisy.
		\item \textbf{Gyroscope} measures angular velocity (\si{\degree\per\second}). Integrating the rate gives angle, but integration drifts over time because of small biases.
	\end{itemize}
	A complementary filter combines the two: 98\,\% trust in the integrated gyro (fast, drift‑prone) and 2\,\% trust in the accelerometer (slow, absolute reference). The equation is
	\[
	\theta_k = 0.98\,(\theta_{k-1} + \omega \cdot \Delta t) + 0.02\,\theta_{\text{accel}}
	\]
	where \(\omega\) is the gyro pitch rate and \(\Delta t = \SI{5}{\milli\second}\).
	
	\subsection{PID Controller}
	The controller’s job is to minimise the error \(\text{setpoint} - \text{measured angle}\). The setpoint is \SI{0}{\degree} (upright).  
	The parallel PID law is
	\[
	u(t) = K_p e(t) + K_i \int_0^t e(\tau)\,d\tau + K_d \frac{de(t)}{dt}.
	\]
	In code this is discretised with \(\Delta t = 5\,\text{ms}\), and the integral is clamped to prevent windup.  
	The output \(u\) is a signed number that represents the PWM duty cycle (±1599). Positive value → forward, negative → backward.
	
	\subsection{Motor Driver}
	The \func{motor\_set\_speed} function translates the signed command into two things:
	\begin{itemize}
		\item \textbf{Direction GPIOs} -- set according to the table in the README.
		\item \textbf{PWM duty cycle} -- the absolute value of the command is written to the TIM2 compare register.
	\end{itemize}
	The same command is sent to both motors, so the robot always tries to balance straight. Differential turning is not yet implemented, but the architecture is ready for it.
	
	\subsection{Encoder Monitoring}
	Each motor has a magnetic encoder that produces a pulse every few degrees of rotation. TIM2 captures the time between two rising edges, calculates the frequency in Hz, and stores it.  
	The encoder data is \textbf{not} used for the balancing controller; it is only transmitted in the telemetry stream. This provides a real‑time view of motor behaviour and can later be used for a cascaded speed controller.
	
	\subsection{Telemetry}
	Every 10 control cycles (i.e., \SI{50}{\milli\second}) the ISR sets a flag.  
	The main loop sees the flag and calls \func{send\_telemetry\_nonblocking}, which formats and sends a comma‑separated line:
	\begin{verbatim}
		angle_deg, gyro_rate_dps, motor_cmd, left_encoder_Hz, right_encoder_Hz
	\end{verbatim}
	A busy flag prevents overlapping transmissions, and the UART TX‑complete callback resets the flag.
	
	% ================================================================
	\section{Detailed Code Explanation}
	\label{sec:code}
	% ================================================================
	
	Every source file is now dissected. For each file we first state its purpose, then list the most important global/static variables, and finally walk through every function.
	
	% ------------------------------------------------------------------
	\subsection{File: \file{mpu6050.c} \& \file{mpu6050.h}}
	\label{subsec:mpu6050}
	% ------------------------------------------------------------------
	\paragraph*{What this module does}
	It talks to the MPU‑6050 over I²C, wakes the sensor, sets the measurement ranges, reads raw acceleration and rotation data, converts the numbers into physical units (\si{g} and \si{\degree\per\second}), and applies bias corrections obtained from a calibration routine.
	
	\paragraph*{Key constants (from \file{mpu6050.h})}
	\begin{lstlisting}[caption=MPU‑6050 address and scale factors, label=lst:mpu_const]
		#define MPU6050_ADDR (0x68U << 1)        // 7-bit address shifted for HAL
		#define MPU6050_ACCEL_SCALE 16384.0f     // LSB per g for +/-2g range
		#define MPU6050_GYRO_SCALE 131.0f        // LSB per deg/s for +/-250 deg/s
	\end{lstlisting}
	The I²C address \texttt{0x68} is prescribed by the MPU‑6050 datasheet. The HAL library expects the address left‑shifted by one bit, hence \texttt{(0x68U << 1)}.  
	The scale factors convert the raw 16‑bit integers to engineering units: for the accelerometer, a reading of 16384 corresponds to \SI{1}{g}; for the gyroscope, 131 corresponds to \SI{1}{\degree\per\second}.
	
	\paragraph*{Internal offset storage}
	\begin{lstlisting}[caption=Global offset structure, label=lst:mpu_offsets]
		static mpu6050_offsets_t g_offsets = {0};
	\end{lstlisting}
	This variable holds the biases calculated during calibration. It is \texttt{static}, so only accessible inside \texttt{mpu6050.c}.
	
	\subsubsection*{Function: \func{mpu6050\_init}}
	\begin{lstlisting}[caption=Sensor initialisation, label=lst:mpu_init]
		HAL_StatusTypeDef mpu6050_init(I2C_HandleTypeDef *hi2c) {
			uint8_t value = 0x00U;
			// Wake up the MPU-6050 (clear sleep bit in PWR_MGMT_1)
			HAL_I2C_Mem_Write(hi2c, MPU6050_ADDR, REG_PWR_MGMT_1, I2C_MEMADD_SIZE_8BIT, &value, 1U, 5);
			// Set accelerometer range to +/- 2g (ACCEL_CONFIG = 0x00)
			HAL_I2C_Mem_Write(hi2c, MPU6050_ADDR, REG_ACCEL_CONFIG, I2C_MEMADD_SIZE_8BIT, &value, 1U, 5);
			// Set gyroscope range to +/- 250 deg/s (GYRO_CONFIG = 0x00)
			HAL_I2C_Mem_Write(hi2c, MPU6050_ADDR, REG_GYRO_CONFIG, I2C_MEMADD_SIZE_8BIT, &value, 1U, 5);
			return HAL_OK;
		}
	\end{lstlisting}
	\textbf{What happens:}
	\begin{enumerate}
		\item Writing \texttt{0x00} to register \texttt{0x6B} (PWR\_MGMT\_1) clears bit 6 (SLEEP). The sensor wakes up and starts converting.
		\item Writing \texttt{0x00} to \texttt{0x1C} (ACCEL\_CONFIG) sets the full‑scale range to \texttt{±2g}. The default is already 0, but the write ensures it.
		\item Writing \texttt{0x00} to \texttt{0x1B} (GYRO\_CONFIG) sets the gyro range to \texttt{±250°/s}.
	\end{enumerate}
	All transfers use a \SI{5}{\milli\second} timeout. If any I²C transaction fails, the function returns \texttt{HAL\_ERROR}.
	
	\subsubsection*{Helper: \func{to\_i16}}
	\begin{lstlisting}
		static int16_t to_i16(uint8_t msb, uint8_t lsb) {
			return (int16_t)((uint16_t)msb << 8U | (uint16_t)lsb);
		}
	\end{lstlisting}
	The MPU‑6050 splits each 16‑bit measurement into two 8‑bit registers (high byte first). This helper reconstructs the signed integer.
	
	\subsubsection*{Function: \func{mpu6050\_read\_accel}}
	\begin{lstlisting}
		HAL_StatusTypeDef mpu6050_read_accel(I2C_HandleTypeDef *hi2c, mpu6050_axes_t *accel_out) {
			uint8_t buf[6];
			HAL_I2C_Mem_Read(hi2c, MPU6050_ADDR, REG_ACCEL_XOUT_H, I2C_MEMADD_SIZE_8BIT, buf, 6, 5);
			accel_out->x = (float)to_i16(buf[0], buf[1]) / MPU6050_ACCEL_SCALE;
			accel_out->y = (float)to_i16(buf[2], buf[3]) / MPU6050_ACCEL_SCALE;
			accel_out->z = (float)to_i16(buf[4], buf[5]) / MPU6050_ACCEL_SCALE;
			// subtract calibration offsets
			accel_out->x -= g_offsets.accel.x;
			accel_out->y -= g_offsets.accel.y;
			accel_out->z -= g_offsets.accel.z;
			return HAL_OK;
		}
	\end{lstlisting}
	\textbf{What happens:}
	\begin{itemize}
		\item A burst read starting at register \texttt{0x3B} grabs six bytes: X high/low, Y high/low, Z high/low.
		\item Each pair is reassembled and divided by the scale factor to get acceleration in \si{g}.
		\item The stored offsets are subtracted, removing any constant bias.
	\end{itemize}
	Because the calibration routine adjusts \(a_z\) to read \SI{1.0}{g} when the sensor is flat, this function returns a \(z\)‑component of \(+1\) in that orientation.
	
	\subsubsection*{Function: \func{mpu6050\_read\_gyro}}
	Works identically but reads from \texttt{0x43} (GYRO\_XOUT\_H) and divides by \texttt{MPU6050\_GYRO\_SCALE}. After subtraction of calibration offsets, the output is the angular velocity in \si{\degree\per\second}.
	
	\subsubsection*{Function: \func{mpu6050\_calibrate}}
	\begin{lstlisting}
		HAL_StatusTypeDef mpu6050_calibrate(I2C_HandleTypeDef *hi2c, uint16_t samples) {
			...
			for (uint16_t i = 0U; i < samples; i++) {
				mpu6050_read_accel(hi2c, &accel);   // reads with current offsets (initially zero)
				mpu6050_read_gyro(hi2c, &gyro);
				ax += accel.x; ay += accel.y; az += accel.z;
				gx += gyro.x; gy += gyro.y; gz += gyro.z;
			}
			g_offsets.accel.x = ax / samples;
			g_offsets.accel.y = ay / samples;
			g_offsets.accel.z = (az / samples) - 1.0f;   // remove 1g from Z
			g_offsets.gyro.x = gx / samples;
			...
		}
	\end{lstlisting}
	\textbf{Why subtract 1.0 on Z?}  
	When the robot sits flat, the only acceleration acting on the sensor is gravity. Gravity pulls with \(1g\) along the negative Z‑axis of the sensor, but the accelerometer reports \(+1g\) because it measures the force exerted by the structure. By subtracting \(1\) from the Z average, we make the Z reading exactly \(1g\) after offset removal. This is crucial for the \(\arctan2(a_y, a_z)\) calculation in the angle filter.
	
	\noindent\textbf{Impact on the project:}  
	Without calibration the robot would think it is tilted even when it is standing perfectly upright. The calibration removes steady‑state errors and is run once before the control loop starts.
	
	% ------------------------------------------------------------------
	\subsection{File: \file{complementary.c} \& \file{complementary.h}}
	\label{subsec:complementary}
	% ------------------------------------------------------------------
	\paragraph*{What this module does}
	It stores the current tilt angle and gyro rate, and provides an update function that fuses the raw sensor data using the complementary filter.
	
	\paragraph*{Static variables}
	\begin{itemize}
		\item \texttt{s\_i2c} -- pointer to the I²C handle used for sensor access.
		\item \texttt{s\_angle\_deg} -- the estimated pitch angle (\si{\degree}).
		\item \texttt{s\_gyro\_rate\_dps} -- the raw gyro pitch rate (\si{\degree\per\second}) for telemetry.
	\end{itemize}
	
	\subsubsection*{Function: \func{angle\_init}}
	Resets the angle and rate to zero and saves the I²C handle. It returns \texttt{HAL\_OK}.
	
	\subsubsection*{Function: \func{angle\_update}}
	\begin{lstlisting}
		HAL_StatusTypeDef angle_update(float dt) {
			mpu6050_axes_t accel, gyro;
			mpu6050_read_accel(s_i2c, &accel);
			mpu6050_read_gyro(s_i2c, &gyro);
			s_gyro_rate_dps = gyro.x;                  // pitch rate
			float accel_angle = atan2f(accel.y, accel.z) * 57.2957795f;
			s_angle_deg = 0.98f * (s_angle_deg + (s_gyro_rate_dps * dt))
			+ 0.02f * accel_angle;
			return HAL_OK;
		}
	\end{lstlisting}
	\textbf{Step‑by‑step:}
	\begin{enumerate}
		\item Read accelerometer and gyro. The gyro X‑axis is the robot’s pitch axis.
		\item Compute the accelerometer‑only angle: \(\theta_{\text{accel}} = \text{atan2}(a_y, a_z) \times 57.3\). This gives the tilt angle in degrees relative to the direction of gravity.
		\item First‑order complementary filter:
		\[
		\theta_{\text{new}} = 0.98 \times (\theta_{\text{old}} + \omega \cdot \Delta t) \;+\; 0.02 \times \theta_{\text{accel}}.
		\]
		The integrated gyro rate gives a responsive but drifting estimate; the accelerometer term pulls it back towards the absolute reference. The weights 0.98 and 0.02 define the time constant of the filter.
	\end{enumerate}
	\textbf{Why not use a Kalman filter?}  
	A complementary filter is much simpler, requires no tuning of process/measurement covariances, and works extremely well for a balancing robot because the dominant motion is at low frequencies where the accelerometer is reliable.
	
	\noindent\textbf{Impact:} This function is called inside the timer ISR at exactly \SI{200}{\hertz}, guaranteed by the hardware timer. Every other module depends on the angle it produces.
	
	\subsubsection*{Getters}
	\texttt{angle\_get\_deg()} and \texttt{angle\_get\_gyro\_rate\_dps()} simply return the current values. They are used by the main ISR to update global telemetry variables.
	
	% ------------------------------------------------------------------
	\subsection{File: \file{pid.c} \& \file{pid.h}}
	\label{subsec:pid}
	% ------------------------------------------------------------------
	\paragraph*{What this module does}
	Implements a generic parallel‑form PID controller. It can be reused for any control loop by providing different gains.
	
	\paragraph*{Data structure}
	\begin{lstlisting}
		typedef struct {
			float kp, ki, kd;
			float integral_limit;
			float output_limit;
		} pid_gains_t;
	\end{lstlisting}
	The structure contains the three gains plus anti‑windup limit for the integral and an output clamping limit.
	
	\paragraph*{Static variables}
	\begin{itemize}
		\item \texttt{s\_gains} -- stored copy of the gains.
		\item \texttt{s\_integral} -- accumulated integral error.
		\item \texttt{s\_prev\_error} -- previous error, used to compute the derivative.
	\end{itemize}
	
	\subsubsection*{Function: \func{pid\_init}}
	Copies the provided gains and resets the integral and previous error to zero.
	
	\subsubsection*{Function: \func{pid\_compute}}
	\begin{lstlisting}
		float pid_compute(float setpoint, float measurement, float dt) {
			float error = setpoint - measurement;
			if (dt <= 0.0f) return 0.0f;
			s_integral += error * dt;
			// clamp integral
			if (s_integral > s_gains.integral_limit) s_integral = s_gains.integral_limit;
			if (s_integral < -s_gains.integral_limit) s_integral = -s_gains.integral_limit;
			float derivative = (error - s_prev_error) / dt;
			s_prev_error = error;
			float output = s_gains.kp * error + s_gains.ki * s_integral + s_gains.kd * derivative;
			// clamp output
			if (output > s_gains.output_limit) output = s_gains.output_limit;
			if (output < -s_gains.output_limit) output = -s_gains.output_limit;
			return output;
		}
	\end{lstlisting}
	\textbf{Explanation:}
	\begin{itemize}
		\item \textbf{Error}: how far the robot is from upright.
		\item \textbf{Integral}: sum of past errors, multiplied by \(\Delta t\). It helps eliminate steady‑state drift. The clamping prevents the integral from growing excessively when the robot cannot right itself (e.g., when held).
		\item \textbf{Derivative}: the rate of change of error, which acts as a damping force to reduce oscillations.
		\item \textbf{Output}: the weighted sum of the three terms, limited to the PWM range.
	\end{itemize}
	\textbf{Why clamp the integral?} Without clamping, the integral term could become huge during a long disturbance, causing the motor command to rail and the robot to overshoot violently when it finally returns. The limit (\texttt{integral\_limit} = 60) was chosen to be a fraction of the PWM maximum.
	
	\textbf{Impact:} The PID block is the mathematical brain. Changing the gains alters the robot’s stiffness and recovery behaviour. All tuning happens here.
	
	% ------------------------------------------------------------------
	\subsection{File: \file{motor.c} \& \file{motor.h}}
	\label{subsec:motor}
	% ------------------------------------------------------------------
	\paragraph*{What this module does}
	Starts the PWM outputs and sets the speed and direction of each motor based on a signed 16‑bit command.
	
	\paragraph*{Constant}
	\begin{lstlisting}
		#define MOTOR_PWM_MAX 1599
	\end{lstlisting}
	This equals the auto‑reload value of TIM2. Writing 1599 to the compare register gives 100\,\% duty cycle.
	
	\paragraph*{Static variable}
	\texttt{s\_pwm\_timer} -- handle to TIM2 for PWM output.
	
	\subsubsection*{Function: \func{motor\_init}}
	Starts PWM generation on channel 3 (left motor) and channel 4 (right motor). At this point the duty cycle is zero because the compare values were initialised to 0.
	
	\subsubsection*{Function: \func{motor\_set\_speed}}
	This is the most hardware‑specific function. Let’s look at the critical part:
	\begin{lstlisting}
		if (motor == MOTOR_RIGHT) {
			if (speed > 0) {
				HAL_GPIO_WritePin(GPIOD, GPIO_PIN_8, GPIO_PIN_SET);   // A high
				HAL_GPIO_WritePin(GPIOD, GPIO_PIN_9, GPIO_PIN_RESET); // B low
			} else if (speed < 0) {
				HAL_GPIO_WritePin(GPIOD, GPIO_PIN_8, GPIO_PIN_RESET);
				HAL_GPIO_WritePin(GPIOD, GPIO_PIN_9, GPIO_PIN_SET);
			} else {
				HAL_GPIO_WritePin(GPIOD, GPIO_PIN_8, GPIO_PIN_RESET);
				HAL_GPIO_WritePin(GPIOD, GPIO_PIN_9, GPIO_PIN_RESET);  // coast
			}
		} ...
	\end{lstlisting}
	The H‑bridge requires two pins per motor. The table below (taken from the README) summarises the logic:
	
	\begin{table}[H]
		\centering
		\begin{tabular}{lcccc}
			\toprule
			\textbf{Motor} & \textbf{DirA pin} & \textbf{DirB pin} & \textbf{Forward} & \textbf{Backward} \\
			\midrule
			Right & PD8 (D12) & PD9 (D8)   & PD8=HIGH, PD9=LOW & PD8=LOW, PD9=HIGH \\
			Left  & PD10 (D10) & PD11 (D7) & PD10=HIGH, PD11=LOW & PD10=LOW, PD11=HIGH \\
			\bottomrule
		\end{tabular}
	\end{table}
	
	After setting direction, the absolute value of \texttt{speed} is written to the PWM compare register:
	\begin{lstlisting}
		__HAL_TIM_SET_COMPARE(s_pwm_timer, channel, magnitude);
	\end{lstlisting}
	\textbf{Impact:} The motor driver abstracts the low‑level pin manipulation. The PID controller just calls \texttt{motor\_set\_speed(MOTOR\_LEFT, cmd)} and the driver does the rest.
	
	% ------------------------------------------------------------------
	\subsection{File: \file{encoder.c} \& \file{encoder.h}}
	\label{subsec:encoder}
	% ------------------------------------------------------------------
	\paragraph*{What this module does}
	Measures the frequency of the encoder pulses from each wheel using TIM2 input capture channels. The data is purely for telemetry; the balancing controller does not use it.
	
	\paragraph*{Data type}
	\begin{lstlisting}
		typedef struct {
			uint32_t period_ticks;
			float frequency_hz;
			uint32_t pulse_count;
		} encoder_data_t;
	\end{lstlisting}
	
	\paragraph*{Static variables}
	\begin{itemize}
		\item \texttt{s\_left\_last\_capture}, \texttt{s\_right\_last\_capture} -- previous capture values for period calculation.
		\item \texttt{s\_left}, \texttt{s\_right} -- volatile structures updated in the ISR.
	\end{itemize}
	
	\subsubsection*{Function: \func{encoder\_init}}
	Starts input capture interrupts on TIM2 channel~1 (left encoder) and channel~2 (right encoder).
	
	\subsubsection*{Function: \func{encoder\_handle\_ic\_callback}}
	This is where the actual measurement happens. It is called by the HAL whenever a capture event occurs.
	\begin{lstlisting}
		// Determine timer clock frequency (taking APB1 prescaler doubling into account)
		HAL_RCC_GetClockConfig(&clk_cfg, &flash_latency);
		timer_clk_hz = HAL_RCC_GetPCLK1Freq();
		if (clk_cfg.APB1CLKDivider != RCC_HCLK_DIV1) {
			timer_clk_hz *= 2U;   // timer clock is doubled if APB1 prescaler > 1
		}
		timer_clk_hz /= (uint32_t)(htim->Init.Prescaler + 1U);
	\end{lstlisting}
	On the STM32F3, when the APB1 prescaler is not 1, the timer clock is double the APB1 clock. This calculation correctly obtains the timer input clock.
	
	Then, depending on the channel, the current capture value is read. The period between the current and previous capture is computed:
	\begin{lstlisting}
		period = (now >= s_left_last_capture) 
		? (now - s_left_last_capture)
		: ((htim->Init.Period - s_left_last_capture) + now + 1U);
	\end{lstlisting}
	This handles timer wrap‑around gracefully. The frequency is then calculated as \texttt{timer\_clk\_hz / period}.
	
	\subsubsection*{Functions: \func{encoder\_get\_left\_data} / \func{encoder\_get\_right\_data}}
	These return a snapshot of the encoder data structure. To prevent reading a partially‑updated value (the ISR may fire at any time), they briefly disable global interrupts with \texttt{\_\_disable\_irq()} / \texttt{\_\_enable\_irq()}. This is a simple form of critical section.
	
	\textbf{Impact:} Having encoder feedback is an \textbf{extra feature}. It allows the user to observe whether motor commands actually result in wheel motion. Later, this data could be fed into a speed controller for more advanced balancing or driving commands.
	
	% ------------------------------------------------------------------
	\subsection{File: \file{main.c}}
	\label{subsec:main}
	% ------------------------------------------------------------------
	This is the glue that brings everything together. We will examine the preprocessor defines, global variables, custom functions, and the main loop.
	
	\paragraph*{Preprocessor defines}
	\begin{lstlisting}
		#define CONTROL_LOOP_HZ 200.0f
		#define CONTROL_DT (1.0f / CONTROL_LOOP_HZ)    // 0.005 s
		#define TELEMETRY_DIVIDER 10U                  // send telemetry every 10th control cycle
	\end{lstlisting}
	
	\paragraph*{Global volatile variables}
	\begin{lstlisting}
		volatile uint8_t telemetry_flag = 1;       // set by ISR to request telemetry
		volatile float g_angle_deg = 0.0f;         // latest filtered angle
		volatile float g_gyro_dps = 0.0f;          // latest gyro rate
		volatile int16_t g_motor_cmd = 0;          // current motor command
		volatile uint32_t g_control_ticks = 0U;    // number of control cycles executed
		volatile uint8_t g_uart_busy = 0U;         // flag to prevent overlapping UART transmissions
	\end{lstlisting}
	Using \texttt{volatile} tells the compiler that these variables can change at any time (from an ISR). Without it, the main loop might never see the updated values.
	
	\subsubsection*{Custom \func{myPrintf} function}
	\begin{lstlisting}
		int myPrintf(const char* fmt, ...) {
			char buffer[512];
			va_list args;
			va_start(args, fmt);
			int length = vsnprintf(buffer, sizeof(buffer), fmt, args);
			va_end(args);
			if (length > 0) {
				HAL_UART_Transmit(&huart2, (uint8_t*)buffer, (uint16_t)length, HAL_MAX_DELAY);
				return 1;
			}
			return 0;
		}
	\end{lstlisting}
	This provides a familiar \texttt{printf}‑like interface over UART. It uses \texttt{vsnprintf} to safely format the string into a stack buffer, then transmits it with the blocking HAL call.  
	\textbf{Note:} Because it uses \texttt{HAL\_MAX\_DELAY}, the function will wait until the whole string is sent. This is acceptable for debug messages but not for the main control loop (which runs in an ISR anyway).
	
	\subsubsection*{Function: \func{start\_robot\_modules}}
	This is the master initialisation routine, called once from \texttt{main()}.
	\begin{lstlisting}
		static void start_robot_modules(void) {
			pid_gains_t gains = {
				.kp = 55.0f, .ki = 1.2f, .kd = 2.8f,
				.integral_limit = 60.0f,
				.output_limit = (float)MOTOR_PWM_MAX,
			};
			mpu6050_init(&hi2c2);
			mpu6050_calibrate(&hi2c2, 300U);   // average 300 samples
			angle_init(&hi2c2);
			pid_init(gains);
			encoder_init(&htim2);
			motor_init(&htim2);
			HAL_TIM_Base_Start_IT(&htim6);     // start the 200 Hz heartbeat
		}
	\end{lstlisting}
	\textbf{Order matters:} the MPU must be initialised and calibrated before the angle filter starts, and the motors and encoder are started last. Once TIM6 begins, the ISR immediately starts calling \texttt{angle\_update} and \texttt{pid\_compute}. If any peripheral fails, \texttt{Error\_Handler()} is called, which disables interrupts and loops forever.
	
	\subsubsection*{Function: \func{send\_telemetry\_nonblocking}}
	\begin{lstlisting}
		static void send_telemetry_nonblocking(void) {
			if (g_uart_busy) return;
			encoder_data_t left = encoder_get_left_data();
			encoder_data_t right = encoder_get_right_data();
			myPrintf("%.3f,%.3f,%d,%.2f,%.2f\r\n", g_angle_deg, g_gyro_dps,
			g_motor_cmd, left.frequency_hz, right.frequency_hz);
			g_uart_busy = 1U;
			if (myPrintf(...)) {   // actually this line is a bug; see text
				g_uart_busy = 0U;
			}
		}
	\end{lstlisting}
	\textbf{Intended behaviour:} The function checks \texttt{g\_uart\_busy} to avoid sending a new line while a previous one is still in progress. After transmitting, the busy flag is set, and the UART TX‑complete callback clears it.
	
	\textbf{Actual behaviour (bug):} The code calls \texttt{myPrintf} twice -- once before setting the busy flag and once in an \texttt{if} statement. The first call sends the data (so the busy flag should already be 1), and the second call overwrites the buffer but never clears the flag correctly. In practice the telemetry still works most of the time because UART is fast at 115200 baud, but it is not a true non‑blocking implementation. This is noted as an area for improvement.
	
	\subsubsection*{Interrupt Callbacks}
	\begin{lstlisting}
		void HAL_TIM_IC_CaptureCallback(TIM_HandleTypeDef *htim) {
			encoder_handle_ic_callback(htim);
		}
		void HAL_TIM_PeriodElapsedCallback(TIM_HandleTypeDef *htim) {
			if (htim->Instance != TIM6) return;
			angle_update(CONTROL_DT);
			g_angle_deg = angle_get_deg();
			g_gyro_dps = angle_get_gyro_rate_dps();
			float command = pid_compute(0.0f, g_angle_deg, CONTROL_DT);
			g_motor_cmd = (int16_t)command;
			motor_set_speed(MOTOR_LEFT, g_motor_cmd);
			motor_set_speed(MOTOR_RIGHT, g_motor_cmd);
			g_control_ticks++;
			if ((g_control_ticks % TELEMETRY_DIVIDER) == 0U) {
				telemetry_flag = 1U;
			}
		}
		void HAL_UART_TxCpltCallback(UART_HandleTypeDef *huart) {
			if (huart->Instance == USART2) {
				g_uart_busy = 0U;
			}
		}
	\end{lstlisting}
	\textbf{Important details:}
	\begin{itemize}
		\item The TIM2 capture callback is simply forwarded to the encoder module; no other capture is used.
		\item The TIM6 period‑elapsed callback is the \textbf{heart of the control loop}. It reads the updated angle, runs the PID, and sets the motors. Everything happens with hard real‑time precision.
		\item The telemetry flag is set every 10 cycles, which the main loop later services.
	\end{itemize}
	
	\subsubsection*{Main Loop}
	\begin{lstlisting}
		int main(void) {
			HAL_Init();
			SystemClock_Config();
			MX_GPIO_Init();
			// ... init all peripherals ...
			start_robot_modules();
			while (1) {
				HAL_Delay(100);
				myPrintf("I was here \r\n");
				if (telemetry_flag) {
					telemetry_flag = 0U;
					send_telemetry_nonblocking();
				}
			}
		}
	\end{lstlisting}
	\textbf{Explanation:}
	\begin{itemize}
		\item \texttt{HAL\_Delay(100)} is a crude way to pace the background tasks. It uses the SysTick timer, which is configured to interrupt every 1\,ms. While waiting, the processor can still service TIM6 and other interrupts, so the control loop is unaffected.
		\item The heartbeat message \texttt{"I was here"} proves the main loop is alive. If it stops printing, the program has probably crashed.
		\item The telemetry flag is checked and cleared atomically (in the main loop context). Because the main loop runs much slower than the ISR, there is no race condition on this single flag.
	\end{itemize}
	
	% ================================================================
	\section{Extra Features Beyond Lab 11 Requirements}
	\label{sec:extra}
	% ================================================================
	
	Lab 11 (Open‑Ended Lab) asks for:
	\begin{itemize}
		\item Angle estimation with a complementary filter.
		\item A PI, PD, or PID controller.
		\item Motors that respond correctly to tilt direction.
		\item Real‑time serial transmission of the filtered angle.
		\item Modular, well‑documented code.
	\end{itemize}
	
	Our project goes well beyond these minimums. The following subsections explain each extension.
	
	\subsection{Encoder‑Based Wheel Speed Measurement}
	\textbf{What we did:} Even though the lab does not require motor feedback, we implemented a full input‑capture driver for both motor encoders. The capture ISR measures the period between each encoder pulse and converts it to a frequency in Hertz.  
	\textbf{Why it’s useful:} The encoder data gives a live view of motor behaviour. It can reveal whether a motor is stalling, spinning faster than expected, or not responding to commands. In future work, this data can close a speed‑control loop, enabling precise velocity control or odometry for distance tracking.
	
	\subsection{Comprehensive Telemetry Stream}
	\textbf{What we did:} The serial output includes five fields:
	\begin{center}
		\texttt{angle, gyro rate, motor command, left encoder Hz, right encoder Hz}
	\end{center}
	\textbf{Why it’s useful:} Debugging a balancing robot is much easier when you can plot not only the angle but also the control output and actual wheel speeds. The extra fields expose the internal state of the controller and the motors, which is invaluable for tuning gains.
	
	\subsection{Non‑Blocking Telemetry Infrastructure}
	\textbf{Intended behaviour:} The introduction of a \texttt{g\_uart\_busy} flag and the UART TX‑complete callback shows an awareness of real‑time constraints. The idea is that the main loop never waits for a transmission to finish; it simply posts data and moves on. The callback will clear the flag when the hardware is free again.  
	\textbf{Current limitation:} The helper \texttt{myPrintf} is still blocking. However, the architecture is ready to be upgraded to interrupt‑ or DMA‑driven UART, at which point the telemetry would truly be non‑blocking. This design thinking is a clear step above a simple busy‑wait print.
	
	\subsection{Custom \texttt{printf} Function \texttt{myPrintf}}
	\textbf{What we did:} Instead of scattering multiple \texttt{HAL\_UART\_Transmit} calls with manual string buffers, a variadic function that mirrors standard \texttt{printf} was written. It uses \texttt{vsnprintf} for safe formatting, preventing buffer overflows.  
	\textbf{Why it’s useful:} This utility makes the code more readable and easier to maintain. It is reused for both the heartbeat message and the telemetry lines.
	
	\subsection{MPU‑6050 Calibration Routine}
	\textbf{What we did:} A \texttt{mpu6050\_calibrate} function averages 300 sensor readings to compute bias offsets for the accelerometer and gyroscope.  
	\textbf{Why it’s essential:} The lab manual mentions a complementary filter but does not explicitly require calibration. In practice, every MPU‑6050 has some bias. Without calibration the angle estimate would have a constant offset, causing the robot to lean. By subtracting the biases, we obtain a cleaner signal. The special treatment of the Z‑axis offset ensures the gravity vector is correctly oriented, which is crucial for the tilt calculation.
	
	\subsection{Modular, Reusable Driver Architecture}
	\textbf{What we did:} The code is split into independent \texttt{.c/.h} pairs, each with a clear interface. Internal variables are declared \texttt{static}, and external access is only through dedicated functions.  
	\textbf{Why it matters:} This follows software engineering best practices. Each module can be tested in isolation. If we wanted to replace the MPU‑6050 with a different IMU, we would only need to change \texttt{mpu6050.c}. The PID, motor, and encoder modules know nothing about the sensor, making the system easy to extend.
	
	\subsection{Fixed‑Rate Control Loop in a Timer ISR}
	\textbf{What we did:} The entire balancing computation runs inside the TIM6 ISR at a guaranteed \SI{200}{\hertz}. No other task can delay it.  
	\textbf{Why it’s superior:} In a busy‑loop implementation, the update rate would vary depending on how long other tasks take. A hardware timer ISR guarantees a constant sample time, which is a requirement for accurate digital control. This is a real‑time system best practice that goes beyond the letter of the lab manual.
	
	\subsection{Atomic Encoder Data Reads}
	\textbf{What we did:} The encoder getter functions temporarily disable interrupts while copying the volatile structure:  
	\begin{lstlisting}
		encoder_data_t encoder_get_left_data(void) {
			encoder_data_t data;
			__disable_irq();
			data = s_left;
			__enable_irq();
			return data;
		}
	\end{lstlisting}
	\textbf{Why it matters:} The capture ISR can fire at any moment and update \texttt{s\_left}. Without the critical section, the main loop might read a half‑updated structure, getting mismatched period and frequency values. This careful handling of concurrency shows attention to detail.
	
	\subsection{Architecture Ready for Differential Steering}
	\textbf{What we did:} Although the current controller applies the same command to both motors, the motor driver supports independent left/right commands, and the encoder module provides separate speed measurements.  
	\textbf{Why it’s valuable:} A natural next step is to implement a turning controller that sends different commands to each wheel. The existing code needs no structural changes to support this; only the PID logic in the ISR would be extended.
	
	\subsection{Heartbeat Message for Liveness Monitoring}
	\textbf{What we did:} The main loop prints \texttt{"I was here"} every 100\,ms.  
	\textbf{Why it’s useful:} When the robot is not connected to a debugger, the heartbeat on a serial terminal is the only way to know that the firmware is still running. If the message stops, the program has likely crashed or hung, prompting investigation.
	
	% ================================================================
	\section{Conclusion}
	\label{sec:conclusion}
	% ================================================================
	
	The firmware described in this report implements a complete sensor‑to‑actuator chain for a self‑balancing robot. It reads an IMU, fuses the data with a complementary filter, runs a PID controller at a rock‑solid \SI{200}{\hertz}, and drives two DC motors. The code is modular, well‑commented, and includes every detail you need to understand how and why it works.
	
	By examining each file line‑by‑line, we have covered the low‑level I²C communication, timer configuration for PWM and input capture, sensor calibration, the complementary filter equation, PID tuning and anti‑windup, motor direction control, UART telemetry, and real‑time system design.
	
	The project also introduces several creative extensions beyond the mandatory lab requirements: encoder‑based wheel speed monitoring, a rich telemetry stream, a non‑blocking transmission framework, a custom \texttt{printf} utility, automatic sensor calibration, modular architecture, atomic data access, and a heartbeat message. These additions demonstrate a deeper understanding of embedded systems and provide a solid foundation for further enhancements.
	
	
\end{document}